\section*{How to read this book}

Everything set in the main column below is Rudin, verbatim. You can read it
straight through and never look at an annotation. The annotation layer sits in
two places, both outside his text.

\medskip
\noindent
\textbf{Margin notes.} A small superscript number in Rudin's text --- like
this\kn{The marker is the only mark added to his page, and it is the same for
every tier so his line spacing is undisturbed. The number out here tells you
what kind of note it is.} --- keys to a numbered note in the outer margin. Since
the book prints in black and white, the \emph{tier} of a note is carried by how
its number is set:

\medskip
\begin{tabular}{@{}p{0.11\textwidth}p{0.82\textwidth}@{}}
\tierplain{7} & \textbf{Plain number --- unpacking.} What the sentence is
actually saying, where a formula came from, an omitted algebraic step.\\[5pt]
\tierdeep{7} & \textbf{Outlined number --- going deeper.} Context beyond the
course: what breaks in general, which later theorem this sets up. Skippable.\\[5pt]
\tierwarn{7} & \textbf{Reversed number --- warning.} A specific belief students
form here that is false, or a place where Rudin's wording will mislead you.
\end{tabular}

\medskip
\noindent
\textbf{Interleaved boxes.} When a note runs longer than a margin will hold, it
becomes a box between Rudin's paragraphs, framed by the same three tiers: a
\textbf{solid} left rule for unpacking, a \textbf{dashed} left rule for going
deeper, and rules \textbf{top and bottom} with a black tab for a warning.

\medskip
\noindent
\textbf{Every figure is an addition.} Rudin's Chapter 1 contains no diagrams
at all. Any picture you see here is annotation, drawn to expose geometry that
his algebra states without illustrating; each carries a grey italic caption.

\medskip
\noindent
\textbf{Where the text is corrected.} The main column follows a retyped edition
whose own foreword concedes that fresh typos were probably introduced. Where one
is evident, the reading is restored from Rudin's third edition and a
warning-tier note records the change --- silent emendation would make
``verbatim'' meaningless.

\medskip
\noindent
\textbf{Structural commentary.} Two further boxes carry no mathematics at all.
They explain the book rather than its contents --- why a chapter exists, why
Rudin ordered it as he did, why a section that looks like a detour is not one.

\begin{roadmap}[Heavy frame, black title bar]
Opens and closes every chapter: what it is for, how it is organized and why in
that order, which parts are load-bearing. Where Rudin has explained his own
choices --- usually in the Preface --- these boxes quote him rather than guess.
\end{roadmap}

\begin{signpost}[Thin frame, inline heading]
Opens a section: why this material sits here, and what to hold on to.
\end{signpost}

\medskip
\noindent
\textbf{Books cited in the notes.} Where an annotation says another treatment
does something differently, it means one of these. They are named because
seeing the same theorem set up two ways is often what makes Rudin's choice
legible.

\medskip
\begin{tabular}{@{}p{0.30\textwidth}p{0.63\textwidth}@{}}
\emph{Tao, Analysis I--II}
  & Terence Tao, 3rd ed. Builds $\N,\Z,\Q,\R$ from the Peano axioms and
    constructs $\R$ by Cauchy sequences rather than cuts. Volume II covers
    metric spaces and uniform convergence --- Rudin's Chapters 2, 4 and 7.\\[5pt]
\emph{Yu Pin, Tsinghua notes}
  & \cjk{数学分析讲义}{Shuxue Fenxi Jiangyi}, Tsinghua / Yau
    Mathematical Sciences Center. Axiomatises $\R$
    by field, order, archimedean and \emph{nested interval} axioms, then derives
    the supremum principle.\\[5pt]
\emph{ECNU text}
  & \cjk{数学分析}{Shuxue Fenxi} (Mathematical Analysis), 5th ed., 2019, East
    China Normal University. Standard Chinese two-volume text. Its \S7.1
    collects six equivalent formulations of completeness and proves them
    interderivable.\\[5pt]
\emph{USTC notes}
  & \cjk{中国科学技术大学少年班数学分析讲义}{USTC Special Class for the
    Gifted Young, analysis notes} (\cjk{程艺}{Cheng Yi} et al., 2005).
    Teaches the main thread of calculus first and defers real-number theory to a later volume;
    lists Rudin among its three principal references.
\end{tabular}
